% Copyright (c) 2026 Geovana Neves. Licensed under CC BY 4.0 (https://creativecommons.org/licenses/by/4.0/). See LICENSE-AND-AUTHORSHIP.md.
%
% workspaces/text/03-cluster-and-keeping.tex
%
% Parts 7 and 8 of guide 2: a workstation and a cluster; keeping, correcting
% and cleaning.

\section{A workstation and a cluster}

\begin{frame}{Two machines, two jobs}
\begin{columns}[T]
\begin{column}{0.48\textwidth}
\begin{block}{The workstation}
\begin{itemize}\small
  \item Write the inputs and plan.
  \item Run one point, or a cheap solver, and look at it.
  \item Collect, read the tables, plot.
\end{itemize}
\end{block}
\end{column}
\begin{column}{0.48\textwidth}
\begin{block}{The cluster}
\begin{itemize}\small
  \item Run the long matrix, one scheduler job per point.
  \item Keep the heavy solver output where it was written.
\end{itemize}
\end{block}
\end{column}
\end{columns}
\begin{itemize}
  \item The workspace is the same tree on both. What differs is one file,
        the \textbf{HPC profile}, that says how this site submits a job.
  \item A workspace with a profile submits; \code{--local} keeps the work on
        this machine all the same.
\end{itemize}
\end{frame}

\begin{frame}[fragile]{The HPC profile}
\begin{columns}[T]
\begin{column}{0.5\textwidth}
\begin{lstlisting}[style=cfdbtoml]
# inputs/hpc/h001.toml
application_id = "SU2"
max_walltime = "96:00:00"
[descriptor]
format = "yaml"
name = "submit.yaml"
[descriptor.fields]
job_name = "{job_name}"
walltime = "{walltime}"
ncpus = "{ncpus}"
version = "{solver_version_alias}"
working_directory = "{workdir}"
config_file = "{config_path}"
[submit]
command = ["submit-tool", "{descriptor_path}"]
[builds]
"8.3.0" = "8.3.0"
\end{lstlisting}
\end{column}
\begin{column}{0.46\textwidth}
\begin{itemize}\small
  \item The profile describes a job file and the command that submits it.
        \code{submit-tool} stands for your site's own.
  \item \code{\{ncpus\}} and \code{\{walltime\}} come from the row's
        \code{NCPUS} and \code{WALLTIME}.
  \item \code{[builds]} maps the row's \code{SOLVER_BUILD} to the name the
        site gives that version.
  \item \code{\{config_path\}} is the file the solver reads in the point
        folder.
  \item \code{--hpc NAME} chooses the default; a row's \code{hpc: NAME}
        in \code{VAR_NAMES_VALUES} overrides it.
\end{itemize}
\end{column}
\end{columns}
\end{frame}

\begin{frame}[fragile]{Select a profile and cluster per row}
\begin{lstlisting}
VAR_NAMES_VALUES
hpc: alpha / cluster: west
hpc: beta / cluster: east
\end{lstlisting}
\begin{itemize}\small
  \item \code{hpc: alpha} selects \code{inputs/hpc/alpha.toml} before the
        \code{--hpc} default. A profile and its digest are recorded per point.
  \item \code{cluster: west} supplies the profile's \code{\{cluster\}}
        placeholder; it falls back to \code{cluster} in \code{init.toml}.
        Keys are case-insensitive and never passed to the solver setup.
  \item \code{--local} ignores both. An unused cluster selection warns;
        \code{--warning-refusals} makes that a refusal.
  \item Plan again after changing a profile. Older HPC receipts without the
        selected profile and digest are refused. Collect waits for all
        configured \code{job_end_files} and settled outputs.
\end{itemize}
\end{frame}

\begin{frame}{A cluster run, step by step}
\begin{enumerate}
  \item \code{cfd-matrix plan polar} prints each point, its job file and the
        submit command. Nothing is written but the receipt.
  \item \code{cfd-matrix run polar} stages every file the solver reads into
        the point folder, then submits one job per point. The points are
        recorded \code{SUBMITTED}.
  \item The job runs in its point folder and needs nothing else from the
        workspace.
  \item \code{cfd-matrix collect polar} reads the points whose output files
        exist and leaves the others outstanding. \code{--watch} repeats the
        sweep until none is outstanding.
\end{enumerate}
\begin{cfdbpitfall}[Send the job's output to the point folder]
Point the scheduler's output files at \code{\{stdout_log\}} and
\code{\{stderr_log\}}: several solvers print their results there, and
collect reads them from the point folder.
\end{cfdbpitfall}
\end{frame}

\begin{frame}[fragile]{Bringing results home: sync}
\begin{lstlisting}
cfd-workspace sync runs
cfd-workspace sync post --apply
cfd-workspace sync all --apply
\end{lstlisting}
\begin{itemize}
  \item \code{inputs/sync-workspaces.toml} names the other workspaces;
        \code{--from NAME} takes one of them.
  \item Each level adds to the one before: \code{runs}, then \code{post},
        then \code{all}.
  \item Sync \textbf{previews} until \code{--apply}.
  \item On a conflict it keeps this workspace's record and file, unless
        \code{--prefer-other} (records) or \code{--overwrite} (files, the
        local copy archived first) says otherwise.
\end{itemize}
\begin{takeaway}Bring the records and the tables to the workstation. Leave
the heavy files on the cluster until you need one.\end{takeaway}
\end{frame}

\section{Keeping, correcting, cleaning}

\begin{frame}{Correcting the records}
\relax{\cfdbbodysize
\begin{tabularx}{\linewidth}{@{}P{0.4\linewidth}R@{}}
\toprule
\textbf{Command} & \textbf{Does} \\
\midrule
\code{cfd-workspace mark-failed --sims POL} & records your verdict that points
  failed, whatever the solver reported \\
\code{cfd-workspace mark-converged --sims POL --reason TEXT} & records your
  verdict that completed points converged; the reason is required \\
\code{cfd-workspace restore KIND} & brings \code{runs}, \code{storage},
  \code{products} or \code{plan} back from \code{archive/<stamp>/} \\
\code{cfd-workspace rebuild} & writes records for simulation folders that have
  none \\
\code{cfd-workspace archive SIM_ID} & moves one simulation's folder and
  records to \code{archive/<stamp>/} \\
\bottomrule
\end{tabularx}}
\vspace{2mm}
\begin{itemize}\small
  \item The correction commands preview until \code{--apply}.
  \item \code{cfd-workspace archive} acts at once, without a preview.
  \item Your verdicts are kept in \code{marks.json}, apart from what the
        solver reported.
\end{itemize}
\end{frame}

\begin{frame}{Space on disk}
\relax{\cfdbbodysize
\begin{tabularx}{\linewidth}{@{}P{0.4\linewidth}R@{}}
\toprule
\textbf{Command} & \textbf{Does} \\
\midrule
\code{cfd-workspace space-in-use} & sizes by folder, by simulation and by file
  extension \\
\code{cfd-workspace free-space RECIPE} & deletes the files that
  \code{inputs/management/<RECIPE>.toml} names, except those a record needs \\
\code{cfd-workspace delete-sims --sims POL} & deletes simulation folders and
  retires their records \\
\bottomrule
\end{tabularx}}
\vspace{2mm}
\begin{itemize}\small
  \item \code{free-space} and \code{delete-sims} preview until
        \code{--apply}. Read the preview: these are the only commands that
        delete.
  \item Every storage command is recorded in \code{storage_management.json};
        \code{cfd-workspace log --storage} lists them.
\end{itemize}
\end{frame}

\begin{frame}{Version control}
\begin{columns}[T]
\begin{column}{0.48\textwidth}
\begin{block}{Worth committing}
\begin{itemize}\small
  \item \code{init.toml}
  \item \code{inputs/}: small text, and the whole definition of the study
  \item \code{runs.json} and \code{post/}, when the results are part of the
        record you keep
\end{itemize}
\end{block}
\end{column}
\begin{column}{0.48\textwidth}
\begin{block}{Usually left out}
\begin{itemize}\small
  \item \code{sims/} and \code{archive/}: solver output
  \item large meshes: keep them where your site keeps large files
  \item \code{inputs/executables.local.toml}: one machine's paths, ignored
        by the \code{.gitignore} that \code{init} writes
\end{itemize}
\end{block}
\end{column}
\end{columns}
\begin{takeaway}The inputs and the records are enough to say what was run.
The solver output is what it cost to run it.\end{takeaway}
\end{frame}
